% SOURCE_AUTHORITY: sga5_fr_workpass.tex, lines 12046--12070 (LF-normalized unit).
% SOURCE_SHA256: 791F4EFFC5E02832D5D77ED03518C8156D6F07E4C8238B03545DB93D883FBB28
\begin{statement}{Definição}
Diremos que o complexo $X$ de $D(\cC_A)$ tem dimensão Tor finita se estiver em $D^-$ e se existir $n_o\in\mathbb Z^+$ tal que
\[
\operatorname{Tor}_n^A(X,Y)=0
\]
para todo $Y\in\Ob({}_A\cC)$ e todo $n>n_o$.
\end{statement}

Usaremos a noção de dimensão Tor para obter uma caracterização dos complexos perfeitos:

\begin{statement}{Proposição 8.1}
Sejam $A$ um anel (não necessariamente comutativo) e $L^\bullet$ um complexo de $A$-módulos. As condições seguintes são equivalentes:
\[
\text{a) } L^\bullet\text{ é perfeito;}
\]
\[
\text{b) } L^\bullet\text{ é pseudocoerente, de cohomologia limitada e de dimensão Tor finita.}\footnote{cf. (SGA 6 I 5.8.1).}
\]
\end{statement}

Para fazer a demonstração é utilizado o seguinte lema:

\begin{statement}{Lema}
(Bourbaki, Alg. Comm. I, §2, ex. 15). Todo $A$-módulo $E$ que é plano e de apresentação finita é projetivo.
\end{statement}
